% Vista científica exacta materializada desde una rama condicional.
% Fuente: source/demostraciones_primarias/31_06c_sectores_topologicos_no_omision.tex; rama: HMTIncludeCMBlock.
% No modifica el contenido matemático de la rama de origen.
\chapter[Complex structure \texorpdfstring{\(\widetilde J_M^{\,2}=-I\)}{J M squared = -I} and topological sectors]
{Complex structure \texorpdfstring{\(\widetilde J_M^{\,2}=-I\)}{J M squared = -I} on the fixed stratum, dodecaphase involutions, and topological sectors}
\label{chap:sectores-topologicos-no-omision}
\index{topological sector}\index{Fibonacci, ring based}
\index{Majorana modes}\index{Ising category}\index{braid group}

The involution \(C_M\), the based ring associated with \(F_{\rm av}\), the
braid characters, and the pointed categories belong to distinct
algebraic strata. They are presented with their domains and identification
controls so that a nominal coincidence does not replace the morphism that
would have to relate them. The atlas census, the infinite spectral hierarchy
whose eight historical heights are initial witnesses, and the statistical
completions preserve the proofs of
\cref{thm:censo-atlas-81,thm:descenso-necesario-fav,%
thm:principio-exclusion-pauli,thm:dos-cuatro-pi-estadistica}.

Each result directly declares its domain, its construction, and its
conclusion; no nominal coincidence modifies those dependencies.

\section{Periodic stratum of the dodecaphasic involution}
\label{sec:involucion-cm-censo}

The domain \(\mathscr T_{12}^{\rm word}\), the conjugation
\(C_M\mathbf t=-\operatorname{rev}(\mathbf t)\), its involutive character, and the
census \(\lvert\operatorname{Fix}(C_M)\rvert=3^6=729\) have their primary residence
in the
\cref{def:palabras-cm,thm:involucion-cm-censo}, within the particle band.
Only the periodic condition required for the topological stratum is added here. If \(r\) denotes the cyclic shift
\((r\mathbf t)_m=t_{m-1}\), we set
\[
 \operatorname{Per}_6
 :=\{\mathbf t\in\mathscr T_{12}^{\rm word}:r^6\mathbf t=\mathbf t\}.
\]

\begin{proposition}[Census of the periodic stratum]
\label{prop:estrato-periodico-cm-rev7}
Upon imposing periodicity with period dividing six,
\[
 \bigl|\operatorname{Fix}(C_M)\cap\operatorname{Per}_6\bigr|
 =3^3=27.
\]
This identity is exact.
\end{proposition}

\begin{proof}
The \cref{thm:involucion-cm-censo} gives the general form of the fixed points.
If, moreover, \(r^6\mathbf t=\mathbf t\), the second half of the word coincides
with the first. Writing the first six coordinates, the two
conditions are equivalent to
\[
 (t_0,t_1,t_2,t_3,t_4,t_5)
 =(a,b,c,-c,-b,-a).
\]
The parameters \(a,b,c\) are independent and take three values; therefore,
the intersection has \(3^3\) elements.
\end{proof}

The general example of a fixed point, prior to requiring period six, is
\[
 (a,b,c,d,e,f,-f,-e,-d,-c,-b,-a).
\]
This coordinated form makes it possible to check by hand both the involution and the
count and avoids interpreting the integer \(729\) by means of another homonymous object.

\begin{remark}[Census of the Involution]
\label{rem:procedencia-n40-cm}
The enumeration of the \(3^{12}=531\,441\) words finds no
violation of the involution and obtains
\[
 |\operatorname{Fix}(C_M)|=729,
 \qquad
 |\operatorname{Fix}(C_M)\cap\operatorname{Per}_6|=27.
\]
The primary proof and \cref{prop:estrato-periodico-cm-rev7} explain
both numbers independently of the traversal: the fixed points have six
free coordinates, and the additional period-six condition reduces the
choice to three. The exhaustive census confirms that no
exceptional orbits exist outside that description.
\end{remark}

\begin{criterion}[Falsifiers of the census \(C_M\)]
\label{crit:falsadores-cm}
Any of the following facts would refute the result: a word for which
\(C_M^2\mathbf t\ne\mathbf t\); a fixed point that does not have the preceding coordinate
form; an exhaustive total other than \(729\); or a total other than
\(27\) after imposing \(r^6\mathbf t=\mathbf t\). The nominal coincidence of this
\(729\) with another cardinal does not identify their domains.
\end{criterion}

\section{Quadratic structure transported to the fixed sheet}
\label{sec:estructura-cuadratica-estrato-fijo-cm}

We identify the alphabet of signs with the ternary field via
\[
 \{-1,0,+1\}\xrightarrow{\sim}\mathbb F_3,
 \qquad
 -1\longmapsto2,\quad 0\longmapsto0,\quad +1\longmapsto1.
\]
With this identification,
\(\mathscr T_{12}^{\rm word}\cong\mathbb F_3^{12}\), and both
\(C_M=-\operatorname{rev}\) and the shift \(r\) are
\(\mathbb F_3\)-linear operators. Consequently,
\(\operatorname{Fix}(C_M)\) and
\(\operatorname{Per}_6=\ker(r^6-I)\) are vector subspaces.

The coordinate form of the fixed points defines the linear isomorphism.
\[
 s_{\rm fix}:\mathbb F_3^6\xrightarrow{\sim}\operatorname{Fix}(C_M),
\]
\[
 s_{\rm fix}(a,b,c,d,e,f)
 =(a,b,c,d,e,f,-f,-e,-d,-c,-b,-a).
\]
Let \(J_W\) be the Paley--Witt operator constructed in
\cref{thm:f9-tres-desde-paley}, with the basis, polarization, and orientation
declared there.

\begin{proposition}[Complex structure \(\widetilde J_M^{\,2}=-I\) on the fixed stratum]
\label{prop:raiz-menos-uno-estrato-fijo-cm}
The transport
\[
\widetilde J_M:=s_{\rm fix}J_Ws_{\rm fix}^{-1}
\in\operatorname{End}_{\mathbb F_3}
   \bigl(\operatorname{Fix}(C_M)\bigr)
\]
satisfies
\[
 \widetilde J_M^{\,2}=-I.
\]
\end{proposition}

\begin{proof}
By conjugation,
\[
 \widetilde J_M^{\,2}
 =s_{\rm fix}J_Ws_{\rm fix}^{-1}
  s_{\rm fix}J_Ws_{\rm fix}^{-1}
 =s_{\rm fix}J_W^2s_{\rm fix}^{-1}
 =-I.
\]
\end{proof}

The result is proved with the chart \(s_{\rm fix}\), the basis, the polarization, and the orientation as explicit starting structure. It does not identify \(C_M\) with \(J_W\): the former is an involution of the word space, and the latter a quadratic structure on its fixed stratum. Nor does it prove that the transport is natural with respect to the entire TPK dynamics.

Let us suppose
\[
 W_{27}
 :=\operatorname{Fix}(C_M)\cap\operatorname{Per}_6.
\]
The inclusion
\[
 W_{27}\hookrightarrow\operatorname{Fix}(C_M)
\]
is canonical, and \cref{prop:estrato-periodico-cm-rev7} proves that
\(\dim_{\mathbb F_3}W_{27}=3\).

\begin{corollary}[Quadratic obstruction in the twenty-seven stratum]
\label{cor:obstruccion-cuadratica-w27}
The subspace \(W_{27}\) cannot be stable under
\(\widetilde J_M\).
\end{corollary}

\begin{proof}
If it were stable, the relation
\(\widetilde J_M^2=-I\) would turn it into a module over
\[
 \mathbb F_3[X]/(X^2+1)\cong\mathbb F_9.
\]
Every vector space over \(\mathbb F_9\) has even dimension when considered
over \(\mathbb F_3\), in contradiction with
\(\dim_{\mathbb F_3}W_{27}=3\).
\end{proof}

The cardinal equality between \(W_{27}\) and the twenty-seven lines of Schläfli
does not furnish a map between the two objects. Such a relation would require
the construction of an incidence map preserving the central orientation and the
action of the stabilizer, together with a compatible dynamical closure. The atlas
of \(81\) cells is not interposed between \(W_{27}\) and the fixed stratum of
\(729\) states: the particle atlas belongs to another domain.

\section[Majorana Levels]{Mathematical stratification of the Majorana object into four levels}
\label{sec:cuatro-niveles-majorana}

The inversion \(C_M\), exterior parity, a relativistic Majorana
fermion, and a Majorana zero mode are objects of different types:

\begin{enumerate}
 \item \(C_M=-\operatorname{rev}\) acts on
 \(\mathscr T_{12}^{\rm word}\). Its involutivity and its census are proved
 exactly; this construction does not, by itself, provide a spinor, a
 field equation, or a mass term.

 \item The sign character \(q=-1\) and the exterior completion belong
 to the exchange algebra. Relative to CAR and the declared grading, they
 produce \(N_s^2=N_s\), spectrum \(\{0,1\}\), and the orientation return
 \(2\pi/4\pi\). These identities follow from CAR and the grading; they do not
 follow from the
 mere census of \(C_M\).

 \item A relativistic Majorana fermion is a spinorial object endowed with a
 relativistic representation, compatible charge conjugation, a
 self-conjugacy condition, a mass term, and dynamics. Those data are not part of
 the definition of \(C_M\) and do not follow from its census.

 \item A Majorana zero mode, or equivalently the nonabelian sector of
 an Ising theory in the topological sense, requires simple objects
 \(\mathbf 1,\psi,\sigma\) with
 \[
  \psi\otimes\psi\simeq\mathbf 1,\qquad
  \psi\otimes\sigma\simeq\sigma,\qquad
  \sigma\otimes\sigma\simeq\mathbf 1\oplus\psi,\qquad
  d_\sigma=\sqrt2,
 \]
 together with coherent associator and braiding. The gap, the splitting, and the
 platform belong to an additional physical realization; the preceding
 fusion rules do not define a universal pole mass.
\end{enumerate}

The separation concerns domains rather than nomenclature: word involution,
CAR parity, self-conjugate spinor, and nonabelian simple object are four distinct
mathematical objects. In the absence of an explicit morphism among them, none
is identified merely because it shares the name “Majorana.”.
